

\begin{abstract}
Several 2026 preprints claim improvements of the bound $H_1\le246$ of
Polymath8b for $H_1=\liminf(p_{n+1}-p_n)$. Each combines equidistribution
estimates for primes in arithmetic progressions (new, or adapted from
Stadlmann's) with a \emph{variational step}:
a test function $F$ on a capped region of $\mathbb R^k$ whose sieve ratio
exceeds a threshold. We recompute the variational step of three of these
claims in exact rational arithmetic, with independent programs, and do not
address their analytic input. Ratios are given in the Polymath scaling, where
the threshold is $4$. (i) For the support and parameters of Stadlmann
($k=49$, claimed $H_1\le240$) we certify the capped ratio $4.0120389\ldots>4$
with a test function of our own, and $4.0124198\ldots$ after optimising it on
the capped region; two independent programs agree exactly.
(ii) For the support shared by Kintali and by Song and Yue ($k=48$, claimed
$H_1\le236$) we recompute the published witness of Song and Yue: its cap-free
mass and its capped numerator equal the paper's reference integers ($655$ and
$715$ digits) exactly, so their criterion holds: in their normalisation (threshold
$1$) $48J_T/I_0=1.0006649215907\ldots>2001/2000$. The authors list the regeneration of the
capped numerator as pending. The same computation reproduces Kintali's
printed capped and cap-free ratios to $12$ significant digits, and a test
function of our own gives $4.0060814\ldots>4$ on the same support. The capped
numerator of the witness and both functionals for our own test function are
confirmed by a second, independent exact program.
(iii) For the draft of Charton et al.\ ($k=45$, claimed $H_1\le212$) the
criterion holds for the caps of its Table~3 ($4.0027351\ldots>4$, two
programs), whereas for the different caps stated in its Lemma~9.1 no
polynomial in an $846$-dimensional space of the size the authors use reaches
$3.99$ (ball arithmetic). This locates an inconsistency in the draft; it is
not a refutation. None of this verifies the new equidistribution estimates,
on which every one of these bounds rests, and none of it bears on the twin
prime conjecture, which remains open.
\end{abstract}

\maketitle

